\documentclass[11pt,a4paper]{article}
% Common preamble for the hint sheets (\input from each file)
\usepackage[utf8]{inputenc}
\usepackage[T1]{fontenc}
\usepackage{lmodern}
\usepackage{amsmath,amssymb,amsthm,mathtools}
\usepackage[a4paper,margin=2.3cm]{geometry}
\usepackage{enumitem}
\usepackage{xcolor}
\usepackage[hidelinks]{hyperref}
\renewcommand{\leq}{\leqslant}
\renewcommand{\le}{\leqslant}
\renewcommand{\geq}{\geqslant}
\renewcommand{\ge}{\geqslant}
\setlength{\parindent}{0pt}
\setlength{\parskip}{0.35em}
% \hintstitle{Sheet title}{Class work / Homework}
\newcommand{\hintstitle}[2]{\begin{center}{\LARGE\bfseries Hints: #1\par}\vspace{0.5em}{\large #2\par}\end{center}
\noindent\textit{The hints to each problem are cumulative. Read one, try again, and open the next only if you are still stuck.}\par\medskip\hrule\medskip}
\newcommand{\problem}[1]{\par\bigskip{\large\bfseries Problem #1}\par\nopagebreak\smallskip}
\newcommand{\subpart}[1]{\par\smallskip\textbf{(#1)}\par\nopagebreak}
\newcommand{\hint}[1]{\par\smallskip\noindent\textcolor{black!70}{\textbf{Hint #1.}}\enspace}
\newcommand{\note}[1]{\par\smallskip\noindent\textit{Note.} #1\par}
\providecommand{\lt}{<}
\providecommand{\gt}{>}
\providecommand{\tightlist}{\setlength{\itemsep}{0pt}\setlength{\parskip}{0pt}}
\emergencystretch=3em

\begin{document}
\hintstitle{Functional sequences, series and spaces -- 2. Series of functions}{Class work}

\problem{1}
\hint{1}
All three tests are statements about the tails. Weierstrass: \(\bigl\lvert\sum_{n=N}^M u_n(x)\bigr\rvert\le\sum_{n=N}^M c_n\) and the uniform Cauchy criterion. Dirichlet and Abel: apply Abel summation \(\sum_{n=N}^M u_nv_n=u_MV_M-u_NV_{N-1}-\sum_{n=N}^{M-1}(u_{n+1}-u_n)V_n\) (with \(V_n\) the partial sums of \(v\), or the tails in Abel's case) and estimate everything uniformly in \(x\); monotonicity makes \(\sum\lvert u_{n+1}-u_n\rvert\) telescope.
\hint{2}
\(\sum\frac{\sin nx}{n^\alpha}\): for \(\alpha\gt1\) uniform on \(\mathbb R\) (Weierstrass). For \(0\lt\alpha\le1\) use \(\bigl\lvert\sum_{n=1}^N\sin nx\bigr\rvert\le\frac1{\lvert\sin(x/2)\rvert}\): Dirichlet gives uniform convergence on every \([\delta,2\pi-\delta]\) (and its translates by \(2\pi k\)), convergence everywhere. Near \(0\) it is not uniform: at \(x=\frac1{2N}\) the block \(\sum_{n=N}^{2N}\) has all terms at least \(\frac{\sin(1/2)}{(2N)^\alpha}\), so it is at least \(c\,N^{1-\alpha}\ge c\).
\hint{3}
\(\sum\frac{(-1)^n}{n+x^2}\): the Leibniz tail estimate \(\le\frac1{N+1+x^2}\le\frac1{N+1}\) gives uniform convergence on all of \(\mathbb R\), while the convergence is absolute nowhere. \(\sum x^n(1-x)\) converges on \((-1,1]\) (a geometric series; the terms vanish at \(x=1\)); its tail from \(N\) is \(x^N\) for \(\lvert x\rvert\lt1\), so the convergence is uniform on \([-a,b]\) for \(0\lt a,b\lt1\) but not on \([0,1)\), \([0,1]\) or near \(-1\); on \([0,1]\) the sum is discontinuous at \(1\).

\problem{2}
\hint{1}
For smoothness, bound all derivatives of \(\frac1{n^2+x^2}\) uniformly in \(x\). A clean way: \(\frac1{n^2+x^2}=\frac1{2in}\Bigl(\frac1{x-in}-\frac1{x+in}\Bigr)\), so the \(k\)-th derivative is at most \(\frac{k!}{n^{k+2}}\) in absolute value; apply the M-test to every differentiated series and Theorem 3.
\hint{2}
Asymptotics: \(t\mapsto\frac1{t^2+x^2}\) is decreasing, so \(\int_1^\infty\frac{dt}{t^2+x^2}\le S(x)\le\int_0^\infty\frac{dt}{t^2+x^2}=\frac{\pi}{2x}\); hence \(S(x)\sim\frac{\pi}{2x}\) as \(x\to\infty\) (in fact \(S(x)=\frac{\pi x\coth\pi x-1}{2x^2}\), if you know it).
\hint{3}
\(\sum\frac{\sin nx}n\) converges everywhere by Dirichlet. Its sum is \(\frac{\pi-x}2\) on \((0,2\pi)\) (integrate \(\frac12+\sum_{n=1}^N\cos nt=\frac{\sin(N+\frac12)t}{2\sin(t/2)}\) from \(x\) to \(\pi\) and let \(N\to\infty\) using the Riemann--Lebesgue lemma), and \(0\) at \(x=0\). The jump at \(0\) shows the convergence cannot be uniform on any neighbourhood of \(0\): a uniform limit of continuous functions is continuous.

\problem{3}
\hint{1}
Uniform convergence is the M-test with \(c_n=1/n^2\). Each term \(\{nx\}/n^2\) is right-continuous and is continuous except at the points with \(nx\in\mathbb Z\), where it drops by \(1/n^2\).
\hint{2}
Because the convergence is uniform, one-sided limits can be taken term by term. At an irrational \(x\) every term is continuous, so \(R\) is continuous. At \(x=p/q\) in lowest terms, the terms that jump are those with \(q\mid n\), and \(R(x^-)-R(x)=\sum_{q\mid n}\frac1{n^2}=\frac{\pi^2}{6q^2}\gt0\).
\hint{3}
Integrability: each partial sum is a bounded function with finitely many discontinuities, hence integrable, and a uniform limit of integrable functions is integrable. (Alternatively, use Lebesgue's criterion: the discontinuity set is countable.)

\problem{4}
\hint{1}
Guess the limit: a fixed point satisfies \(f'=\frac1{1+f}\), \(f(0)=0\), i.e.~\(\bigl((1+f)^2\bigr)'=2\), so \(f(x)=\sqrt{1+2x}-1\).
\hint{2}
Write \(T g(x)=\int_0^x\frac{dt}{1+g(t)}\). For \(g,h\ge0\), \(\lvert Tg-Th\rvert(x)\le\int_0^x\lvert g-h\rvert(t)\,dt\), because \(\Bigl\lvert\frac1{1+g}-\frac1{1+h}\Bigr\rvert=\frac{\lvert g-h\rvert}{(1+g)(1+h)}\le\lvert g-h\rvert\). All \(f_n\) are \(\ge0\).
\hint{3}
Iterate this Volterra-type estimate against the fixed point \(f=Tf\): \(\lvert f_n-f\rvert(x)\le\frac{x^n}{n!}\max\lvert f_0-f\rvert\). So \(f_n\to\sqrt{1+2x}-1\), and the convergence is uniform on \([0,1]\).

\problem{5}
\hint{1}
Each \(f_n\) is a polynomial with \(f_n'=f_{n-1}\) and \(f_n(a_n)=0\). Compute the first cases: \(f_1(x)=x-a_1\), and \(f_2(0)=0\) forces \(a_2(a_2-2a_1)=0\).
\hint{2}
Since \(f_n'=f_{n-1}\), every \(f_n\) is \(\frac{x^n}{n!}+c_1\frac{x^{n-1}}{(n-1)!}+\dots+c_n\) with \(c_k=f_k(0)\) (a generating function \(\sum f_n(x)s^n=C(s)e^{sx}\) makes this transparent). The condition says \(c_k=0\) for \(k\ge2\), so by induction \(f_n(x)=\frac{x^n}{n!}-a_1\frac{x^{n-1}}{(n-1)!}=\frac{x^{n-1}(x-na_1)}{n!}\).
\hint{3}
In the induction step, \(f_n\) is that polynomial plus the constant that makes it vanish at \(a_n\); the constant is zero iff \(a_n^{n-1}(a_n-na_1)=0\). Answer: \(a_1\) arbitrary, and for each \(n\ge2\) independently \(a_n=0\) or \(a_n=na_1\).
\end{document}
